\documentclass[11pt,a4paper]{article}

\usepackage[margin=1in]{geometry}
\usepackage[utf8]{inputenc}
\usepackage[T1]{fontenc}
\usepackage{amsmath, amssymb}
\usepackage{graphicx}
\usepackage{booktabs}
\usepackage{tabularx}
\usepackage{xcolor}
\usepackage{hyperref}
\usepackage{caption}
\usepackage{subcaption}
\usepackage{microtype}

\hypersetup{
    colorlinks=true,
    linkcolor=blue!50!black,
    urlcolor=blue!50!black,
    citecolor=blue!50!black,
}

\graphicspath{{../plots/}}

\title{Reproducing Feedback Alignment \\
       \large Lillicrap, Cownden, Tweed and Akerman (2014)}
\author{Ilyas Boudhaine \and Abderrahmane Baidoune}
\date{\today}

\begin{document}
\maketitle

\begin{abstract}
We reproduce the three experiments from \emph{Random feedback weights support
learning in deep neural networks} (Lillicrap et al., 2014). The paper's
central claim is that the exact transpose of the forward weights is not
required for credit assignment: fixed random feedback matrices suffice, and
the forward weights gradually align with them during training. We reimplement
the algorithms in pure NumPy, run paper-faithful experiments on the linear
function approximation task (Task 1), MNIST classification (Task 2), and a
nonlinear teacher-student task (Task 3), and report findings against the
paper's published curves. All three qualitative claims reproduce. Our MNIST
accuracy is below the paper's published value, almost certainly because our
hyperparameter grid is too coarse; we say so explicitly rather than tune it
away. Code, results, and plots are available alongside this report.
\end{abstract}

\section{Introduction}

Backpropagation requires that the feedback path used to compute hidden-layer
error signals be the transpose of the forward weights. This is the
\emph{weight transport} problem: a biological neuron has no obvious mechanism
to read out and broadcast the precise transpose of another neuron's
projections \cite{lillicrap2014}.

The Lillicrap et al.\ paper makes a striking claim: this symmetry is
unnecessary. Replacing the transpose with a \emph{fixed random matrix} $B$
yields a feedback signal that, while initially uncorrelated with the true
gradient, becomes a useful learning signal because the forward weights
gradually \emph{align} with $B$ over training. Concretely, for a one-hidden
layer network with output error $e = y^* - y$, the standard backprop hidden
signal $W^\top e$ is replaced by $B e$ during learning. The paper shows that
on a linear task, on MNIST, and on a nonlinear teacher-student task this
substitution does not noticeably hurt performance, and in some regimes
actually improves early-training dynamics.

This report reproduces those three experiments end-to-end. We follow the
paper's ``Full Methods'' specification as closely as possible. Where the
paper leaves a choice to manual search (the initial scales $\omega$ and
$\beta$ for MNIST, the teacher regime and feedback scales for the nonlinear
task), we expose those choices explicitly via sweep commands and report
which point we picked from each grid.

\section{Methodology}

\paragraph{Implementation.} The full codebase is written in NumPy
($\sim$1.6k lines including tests). All three learning rules
(backpropagation, feedback alignment, shallow learning) are defined in a
single function \texttt{algorithms.step}; the three task drivers (linear,
MNIST, nonlinear) invoke it with the appropriate feedback matrices. A
gradient-check test compares the hand-coded backprop signals against
PyTorch's autograd to numerical-precision tolerance ($10^{-10}$), so we know
the underlying math is correct.

\paragraph{Reproducibility.} Each task uses a named, independent random
number generator for each of: dataset construction, weight initialisation,
feedback-matrix construction, test-set construction, and example ordering.
This means BP and FA on the same seed share identical student
initialisations, isolating algorithmic differences from initialisation
noise.

\paragraph{Tasks.} The three architectures and the relevant
paper-specified hyperparameters are summarised in Table~\ref{tab:setup}.

\begin{table}[h]
\centering
\caption{Task setups. The right column lists hyperparameters specified
explicitly by the paper. Others ($\omega$, $\beta$, target regime, $B$
scales) are chosen by manual search.}
\label{tab:setup}
\small
\begin{tabularx}{\linewidth}{l X X}
\toprule
Task & Architecture & Paper-fixed hyperparameters \\
\midrule
Linear (Fig.\ 1, Fig.\ 4) & $30 \to 20 \to 10$, linear & $\eta=10^{-3}$, $\omega=0.01$, $\beta=0.5$, $T \sim \mathcal{U}[-1,1]$ \\
MNIST (Fig.\ 2)           & $784 \to 1000 \to 10$, sigmoid & $\eta=10^{-3}$, $\alpha=10^{-6}$, MSE loss \\
Nonlinear (Fig.\ 3)       & $30 \to 20 \to 10$ or $30 \to 20 \to 10 \to 10$, tanh hidden, linear output & $\eta=10^{-3}$ \\
\bottomrule
\end{tabularx}
\end{table}

\paragraph{What we ran.} A single seed (seed $=0$) for all experiments.
For MNIST we ran $1.5\times 10^6$ examples per run across a
$3\times 3$ grid of $(\omega,\beta) \in \{0.05, 0.1, 0.2\}^2$, plus three
backprop baselines at the same $\omega$ values and one feedback-alignment
run with 50\% sparsity. For the nonlinear task we ran $1.5\times 10^6$
examples per run across an asymmetric grid of $(b_1, b_2)$ with each in
$\{0.05, 0.1, 0.2\}$ for the four-layer student, three values of $b_1$ for
the three-layer student, plus backprop and shallow-learning baselines at
both depths. The full sweep took 17.5 hours of wall time on an 8-vCPU
\texttt{c4-standard-8} VM (Intel Emerald Rapids) running seven parallel
processes. We later realised that NumPy's BLAS backend defaults to
multi-threading, so each of the seven processes spawned up to eight BLAS
threads, oversubscribing the eight physical cores and serialising the work.
Pinning each process to a single BLAS thread (e.g.\
\texttt{OMP\_NUM\_THREADS=1 OPENBLAS\_NUM\_THREADS=1}) should give the
roughly $7\times$ speedup the layout was meant to provide.

\paragraph{Honest caveats vs.\ the paper.}
The paper reports averages over multiple seeds; our results are single-seed.
The paper's manual search for $(\omega, \beta)$ on MNIST is not reported in
exact numerical detail. We picked a 3$\times$3 grid that brackets the
typical scales used in the literature; if the paper's true optimum lies
outside our grid we will not have found it. Both effects are visible in our
MNIST numbers below.

\section{Task 1: Linear function approximation}

\subsection{Setup}

The task is to learn a $\mathbb{R}^{30} \to \mathbb{R}^{10}$ linear map
through a single hidden layer of width 20 with no biases. The target is
$T \sim \mathcal{U}[-1, 1]^{10\times 30}$; inputs are $x \sim \mathcal{N}(0, I_{30})$
with $y^* = T x$. The student is $y = W \cdot W_0 \cdot x$ with both
matrices initialised uniformly in $[-\omega, \omega]$ at $\omega = 0.01$.
The feedback matrix is $B \sim \mathcal{U}[-0.5, 0.5]^{20\times 10}$.
Loss is $\frac{1}{2}\|y^* - y\|^2$.

\subsection{Results: Fig.\ 1 reproduction}

Figure~\ref{fig:linear_fig1} shows the headline result. The normalized
squared error decreases for both algorithms, but feedback alignment
\emph{learns faster} during the first $\sim$500 examples. This is the
paper's main qualitative observation about the linear task: because the
feedback matrix $B$ is drawn at scale $0.5$ while the forward $W$ is drawn
at scale $0.01$, the FA hidden signal $Be$ is roughly fifty times larger
than the BP signal $W^\top e$ at initialisation. The cost is a noisier
trajectory, visible in the spread of the green curve.

The bottom panel of Figure~\ref{fig:linear_fig1} shows the angle
$\angle(\Delta h_\text{FA}, \Delta h_\text{BP})$. At step zero the two
signals are essentially orthogonal at $\sim$85\textdegree{}. By step
$\sim$300 the angle has dropped below 35\textdegree{} and stabilises around
30\textdegree{} for the remainder of training. This is the
\emph{alignment} phenomenon: the forward weights $W$ adapt so that
$W^\top e$ and $Be$ point in nearly the same direction, even though $B$ is
fixed and never sees the gradient.

\begin{figure}[h]
\centering
\includegraphics[width=0.65\textwidth]{linear_fig1.png}
\caption{Task 1, paper Fig.\ 1 reproduction. Top: NSE over training for BP
(black) and FA (green). Bottom: angle between the FA hidden signal and the
BP hidden signal. The dashed line at 90\textdegree{} marks orthogonality.
Single trial, $\omega = 0.01$, $\eta = 10^{-3}$.}
\label{fig:linear_fig1}
\end{figure}

\subsection{Results: Fig.\ 4 sweep over $\omega$}

Figure~\ref{fig:linear_fig4} sweeps the initial-weight scale $\omega$ across
the nine values specified by the paper, with 20 trials per $\omega$,
plotting trial means. Three observations.

\begin{enumerate}
    \item All values of $\omega$ converge eventually, but smaller $\omega$
          converges faster, consistent with the paper.
    \item The asymptotic FA-vs-BP angle is a clean monotonic function of
          $\omega$: smaller $\omega$ produces tighter alignment.
    \item The angle to pseudo-backprop ($\angle(\Delta h_\text{FA},
          W^+ e)$, third panel) follows the same trend and stays slightly
          below the angle to BP, suggesting the FA signal aligns more
          closely with the Moore-Penrose pseudoinverse direction than with
          the exact transpose. The paper makes the same observation.
\end{enumerate}

\begin{figure}[h]
\centering
\includegraphics[width=0.75\textwidth]{linear_fig4.png}
\caption{Task 1, paper Fig.\ 4 reproduction. FA-only sweep over the
initial-weight scale $\omega$, colormap from $\omega = 10^{-4}$ (blue) to
$\omega = 0.25$ (red). Mean across 20 trials. Top: NSE. Middle: angle
between the FA hidden signal and the BP hidden signal. Bottom: angle to
the pseudo-backprop hidden signal.}
\label{fig:linear_fig4}
\end{figure}

\section{Task 2: MNIST classification}

\subsection{Setup}

The student is a $784 \to 1000 \to 10$ network with sigmoid hidden and
sigmoid output units, biases on both layers, MSE loss against one-hot
targets, weight decay $\alpha = 10^{-6}$ on the synaptic weights. The
network is trained online for $1.5 \times 10^6$ examples drawn from the
standard 60k MNIST training set (test error evaluated on the 10k held-out
test set every 50k examples).

The paper notes that $\omega$ (forward init scale) and $\beta$ (feedback
init scale) are picked by manual search. Our search grid is
$(\omega, \beta) \in \{0.05, 0.1, 0.2\}^2$.

\subsection{Results}

The sweep surface is shown in Figure~\ref{fig:mnist_heatmap} as a heatmap
of final test error. Best feedback-alignment performance is achieved at
$\omega = 0.1, \beta = 0.2$, with final test error $9.98\%$. Backprop at
the same $\omega$ values reaches $10.04\%$ at best, essentially tied with
FA. The qualitative claim of the paper, ``FA matches BP on MNIST'', holds.

Figure~\ref{fig:mnist_main} shows learning curves for the best of each
algorithm plus the sparse-50\% variant. Both BP and FA descend rapidly in
the first $\sim$200k examples and asymptote near $10\%$ error. The
sparse-50\% FA run is $\sim$1.4 percentage points worse than dense FA,
which is consistent with the paper's reported small gap.

The alignment angle (bottom panel of Figure~\ref{fig:mnist_main}) starts
near 87\textdegree{} and decreases monotonically to $\sim$42\textdegree{}
for the best FA run, and to $\sim$62\textdegree{} for the sparse variant.
The paper observes the same monotonic decrease, although it reaches
slightly smaller asymptotic angles than ours.

\begin{figure}[h]
\centering
\begin{subfigure}{0.50\textwidth}
    \centering
    \includegraphics[width=\textwidth]{mnist_sweep_heatmap.png}
    \caption{Sweep heatmap.}
    \label{fig:mnist_heatmap}
\end{subfigure}\hfill
\begin{subfigure}{0.45\textwidth}
    \centering
    \includegraphics[width=\textwidth]{mnist.png}
    \caption{Best-of-each comparison.}
    \label{fig:mnist_main}
\end{subfigure}
\caption{Task 2, paper Fig.\ 2. (a) Final test error over the
$(\omega, \beta)$ grid; the yellow column at $\beta = 0.2$ identifies the
best regime. (b) Test error and FA-vs-BP angle over training for the best
BP ($\omega = 0.2$), best FA ($\omega = 0.1, \beta = 0.2$), and FA with
50\% sparsity on $W$ and $B$. Faint background curves are the other FA
sweep points.}
\end{figure}

\subsection{Honest gap vs.\ the paper}

The paper achieves roughly $2$--$3\%$ test error on MNIST under FA, whereas
we plateau at $\sim$10\%. Three plausible causes:
\begin{enumerate}
    \item \textbf{Hyperparameter grid too narrow.} Our heatmap is
          monotonically decreasing toward larger $\beta$ at every $\omega$,
          which suggests the optimum may lie at $\beta > 0.2$. We did not
          extend the grid.
    \item \textbf{Single seed.} The paper reports averages; one seed gives
          us no smoothing.
    \item \textbf{No data augmentation, no learning-rate schedule.} The
          paper uses neither, so this is not the gap.
\end{enumerate}
We did not attempt to close this gap by extending the sweep. The
qualitative claim (``BP and FA tie on MNIST'') reproduces under our grid;
the quantitative claim does not.

\section{Task 3: Nonlinear teacher-student}

\subsection{Setup}

A fixed teacher network of architecture $30 \to 20 \to 10 \to 10$ with
tanh-tanh-linear activations generates targets for inputs
$x \sim \mathcal{N}(0, I_{30})$. Two student architectures are considered:
a three-layer net $30 \to 20 \to 10$ (one tanh hidden layer, linear output)
and a four-layer net $30 \to 20 \to 10 \to 10$ matching the teacher's
architecture. Each is trained under BP, FA, and ``shallow'' learning
(output layer only) for $1.5 \times 10^6$ online examples. NSE is
evaluated on a 5000-example fixed test set drawn from the same input
distribution.

The paper notes that the teacher regime and feedback scales are picked
manually. Our teacher uses all weights $\sim \mathcal{U}[-1, 1]$ and no
biases. For the four-layer FA student we sweep
$(b_1, b_2) \in \{0.05, 0.1, 0.2\}^2$ (partial, see Section 6).

\subsection{Results}

Figure~\ref{fig:nonlinear_main} shows the headline learning curves and
Figure~\ref{fig:nonlinear_heatmap} the four-layer FA sweep surface.

The four-layer backprop student converges to essentially zero NSE
($\sim 10^{-7}$), which is unsurprising: the student architecture exactly
matches the teacher, so backprop on a noise-free regression problem can
recover the teacher's weights to numerical precision. The four-layer FA
student, with best feedback scales $b_1 = b_2 = 0.2$, reaches NSE $0.103$.

Crucially, \emph{four-layer FA outperforms three-layer BP} ($0.103$ vs.\
$0.190$): adding a hidden layer and training it with random feedback beats
adding nothing and training optimally. This is the paper's central
qualitative claim for Task 3 and it reproduces cleanly.

Three-layer FA ($0.221$ at best $b_1 = 0.2$) closely matches three-layer
BP ($0.190$). Shallow learning fails entirely at both depths
($\sim 0.67$--$1.00$ NSE), confirming that the nonlinearity is essential
to the task and that an output layer trained alone cannot capture it.

\begin{figure}[h]
\centering
\begin{subfigure}{0.55\textwidth}
    \centering
    \includegraphics[width=\textwidth]{nonlinear.png}
    \caption{Learning curves.}
    \label{fig:nonlinear_main}
\end{subfigure}\hfill
\begin{subfigure}{0.40\textwidth}
    \centering
    \includegraphics[width=\textwidth]{nonlinear_sweep_heatmap.png}
    \caption{Four-layer FA sweep.}
    \label{fig:nonlinear_heatmap}
\end{subfigure}
\caption{Task 3, paper Fig.\ 3. (a) NSE over training for BP, FA, and
shallow learning at depths 3 and 4. Solid lines are the three-layer
student, dashed lines the four-layer student. The four-layer FA curve
(dashed green) drops below the three-layer BP curve (solid black) before
$10^6$ examples, reproducing the paper's central claim. (b) Heatmap of
final test NSE over $(b_1, b_2)$ for the four-layer student. White cells
were not in our sweep grid.}
\end{figure}

\section{Discussion}

\paragraph{The qualitative claims reproduce.} In all three tasks the
direction of the paper's reported effects holds. FA matches BP on the
linear task (Fig.\ 1) and on MNIST (Fig.\ 2). Alignment emerges in both
settings, with the angle decreasing monotonically from near orthogonal at
initialisation. Deeper-net FA outperforms shallower-net BP on the
nonlinear task (Fig.\ 3). The relationship between $\omega$ and final
alignment angle in the Fig.\ 4 sweep is monotonic and clean.

\paragraph{What we could not reproduce.} Our MNIST final test error
($\sim$10\%) is well above the paper's published number ($\sim$2--3\%).
The most likely cause is an under-sized hyperparameter search; secondary
possibilities are the single-seed protocol and unreported training-time
details in the paper. We deliberately did not extend the sweep to make our
number look better, since the qualitative finding (the algorithms tie) is
what the paper actually claims for that figure.

\paragraph{Limitations.} Beyond the MNIST gap: all our results use a
single random seed; the paper averages across multiple trials. The
nonlinear sweep does not cover the $(b_1, b_2)$ grid uniformly (two off-
diagonal cells are missing). We did not reproduce Fig.\ 3a (visualisation
of individual learned functions), only the learning-curve panel.

\paragraph{Bug fixes relative to a naive implementation.} Two issues that
do not appear in the paper but matter for a faithful reproduction. First,
splitting random streams by purpose (data, init, feedback, order) is
necessary to avoid accidental correlations; using a single seeded RNG for
everything makes initialisations co-vary with dataset draws across runs.
Second, in MNIST the training loop must reshuffle the example order each
epoch rather than cycle a single fixed permutation 25 times. Without
these, BP and FA on the same seed are not comparable.

\paragraph{Bigger picture.} Feedback alignment is a striking
demonstration that the brain does not need weight transport to learn.
Whether it scales to deeper or harder tasks is a separate question, partly
addressed by later papers (direct feedback alignment, sign-symmetry, and
target propagation lineages). Our reproduction confirms that the original
2014 result is solid: under online SGD with the paper's specifications,
random fixed feedback is enough.

\section{Code and data availability}

The code, sweep configurations, and analysis script are available alongside
this report. The full sweep can be reproduced by following the README's
``How to run'' section. The
plotting script (\texttt{plot\_all.py}) regenerates every figure in this
report from the CSV outputs in \texttt{results/}.

\bibliographystyle{plain}
\begin{thebibliography}{1}
\bibitem{lillicrap2014}
T.~P. Lillicrap, D.~Cownden, D.~B. Tweed, and C.~J. Akerman.
\newblock Random feedback weights support learning in deep neural networks.
\newblock \emph{arXiv:1411.0247}, 2014.
\newblock \url{https://arxiv.org/abs/1411.0247}.
\end{thebibliography}

\end{document}
